
\part{Imitation learning}

\chapter{Learning from demonstrations}

Imitation learning is a paradigm in machine learning where an agent learns to perform tasks by mimicking expert demonstrations rather than relying solely on explicit reward signals.
This approach is particularly useful in scenarios where designing a reward function is challenging or when expert behavior is readily available.

\section{Types of imitation learning}

\begin{itemize}
    \item \textbf{Behavioral Cloning:} A supervised learning approach where the agent learns a mapping from states to actions based on the expert demonstrations.
        The agent is trained to minimise the difference between its actions and the expert's actions.

    \item \textbf{Inverse Reinforcement Learning (IRL):} The agent infers the underlying reward function that the expert is optimizing.
        Once the reward function is learned, the agent can use reinforcement learning techniques to derive an optimal policy.
\end{itemize}
